\begin{exercise}{102}
Show that:
\begin{enumerate}[label=(\arabic*)]
 \item A retract of a Joyal equivalence is a Joyal equivalence.
 \item The set of monomorphisms which are also Joyal equivalences is saturated.
\end{enumerate}
Prove or disprove that the set of Joyal equivalences is saturated.

\begin{proof}
(1) For every $\infty$-category $\mathcal E$, restriction into $\mathcal E$
turns a retract diagram of maps into a retract diagram of functors
between $\infty$-categories. Passing to $h\mathrm{Cat}_\infty$ preserves
this diagram. A retract of an isomorphism is an isomorphism, so
Definitions~2.1.14 and~2.2.9 prove the assertion.

(2) For a monomorphism $i\colon A\to B$, Proposition~2.2.5 applied to
$\mathcal E\to\Delta^0$ says that
$i^*\colon\mathcal E^B\to\mathcal E^A$ is an isofibration.
Thus, by Definition~2.2.9 and Lemma~2.2.17, $i$ is a Joyal
equivalence if and only if every $i^*$ is a trivial fibration.
For $n\ge0$, put
$q_{\mathcal E,n}\colon\mathcal E^{\Delta^n}\to
\mathcal E^{\partial\Delta^n}$.
By cartesian adjunction, a boundary lifting problem for $i^*$ is
exactly a lifting problem for $i$ against $q_{\mathcal E,n}$:
in both cases a filler is a compatible map $B\times\Delta^n\to\mathcal E$.
Consequently the class in (2) is
\[
 \mathrm{Mono}\cap
 \bigcap_{\mathcal E,\,n\ge0}\chi_L(\{q_{\mathcal E,n}\}),
\]
where $\mathcal E$ ranges over the $\infty$-categories.
Each lifting class is saturated by Lemma~1.3.8, as is
$\mathrm{Mono}$ by Exercise~\ref{ex:59}.
Their intersection is saturated by Remark~1.3.7.

The class of all Joyal equivalences is not saturated.
Let $G$ be the groupoid with two objects and exactly one morphism
between each ordered pair, so $J=N(G)$.
The projection $b\colon J\to\Delta^0$ is a Joyal equivalence:
choosing either object gives an inverse up to natural isomorphism.
Regard $\mathbb Z/2=\{e,t\}$ as a one-object groupoid, and let
$\phi\colon G\to\mathbb Z/2$ send both nonidentity arrows to $t$.
The map $N\phi$ is surjective in every degree: lift a string of group
elements by starting at one object and changing objects exactly at
each occurrence of $t$.
Hence the following square is a degreewise pushout:
\[
\begin{tikzcd}[column sep=large,row sep=large]
 J\arrow[r,"b"]\arrow[d,"N\phi"'] & \Delta^0\arrow[d,equal] \\
 N(\mathbb Z/2)\arrow[r,"q"'] & \Delta^0.
\end{tikzcd}
\]
But $q$ is not a Joyal equivalence: by Exercise~\ref{ex:89}, its
homotopy-category functor would be an equivalence, whereas
$\mathbb Z/2\to *$ is not faithful.
Thus Joyal equivalences are not closed under pushouts.
\end{proof}
\end{exercise}
